\section{Introduction}

Carl Friedrich Gauss revolutionized number theory with his seminal work \textit{Disquisitiones Arithmeticae} (1801), where he systematized the theory of congruences. The notion of modular arithmetic---working with remainders upon division---provides the foundation for modern cryptography, coding theory, and computational number theory.

In this chapter, we explore:
\begin{itemize}
    \item Congruences and residue classes
    \item Modular arithmetic operations
    \item The Chinese Remainder Theorem (Gauss's proof)
    \item Primitive roots and discrete logarithms
    \item Gauss's algorithm for computing discrete logarithms
\end{itemize}

\section{Congruences and Residue Classes}

\begin{definition}[Congruence]
Let $a, b, n \in \Z$ with $n > 0$. We say $a$ is \textit{congruent} to $b$ modulo $n$, written
\[
a \equiv b \pmod{n},
\]
if $n \mid (a - b)$, i.e., $a - b = kn$ for some $k \in \Z$.
\end{definition}

The relation $\equiv \pmod{n}$ is an equivalence relation on $\Z$. The equivalence classes are called \textit{residue classes modulo $n$}. The set of all residue classes is denoted $\Z/n\Z$ or $\Z_n$.

\begin{example}
Modulo $7$, the residue classes are:
\[
[0]_7 = \{\ldots, -7, 0, 7, 14, \ldots\},\quad
[1]_7 = \{\ldots, -6, 1, 8, 15, \ldots\},\quad \ldots,\quad
[6]_7 = \{\ldots, -1, 6, 13, 20, \ldots\}.
\]
We typically represent each class by its least non-negative representative $\{0, 1, 2, 3, 4, 5, 6\}$.
\end{example}

\begin{theorem}[Properties of Congruences]
If $a \equiv b \pmod{n}$ and $c \equiv d \pmod{n}$, then:
\begin{enumerate}
    \item $a + c \equiv b + d \pmod{n}$
    \item $a - c \equiv b - d \pmod{n}$
    \item $ac \equiv bd \pmod{n}$
    \item $a^k \equiv b^k \pmod{n}$ for any $k \in \mathbb{N}$
\end{enumerate}
\end{theorem}

\section{The Chinese Remainder Theorem}

Gauss gave the first complete proof of the Chinese Remainder Theorem in \textit{Disquisitiones Arithmeticae}, Article 36--44.

\begin{theorem}[Chinese Remainder Theorem]
Let $n_1, n_2, \ldots, n_k$ be pairwise coprime positive integers. For any integers $a_1, a_2, \ldots, a_k$, the system of congruences
\[
\begin{cases}
x \equiv a_1 \pmod{n_1} \\
x \equiv a_2 \pmod{n_2} \\
\vdots \\
x \equiv a_k \pmod{n_k}
\end{cases}
\]
has a unique solution modulo $N = n_1 n_2 \cdots n_k$.
\end{theorem}

\begin{proof}[Gauss's Constructive Proof]
Let $N = \prod_{i=1}^k n_i$ and $N_i = N/n_i$. Since $\gcd(N_i, n_i) = 1$, there exists an integer $M_i$ such that $N_i M_i \equiv 1 \pmod{n_i}$. Then
\[
x = \sum_{i=1}^k a_i N_i M_i
\]
satisfies all congruences. Uniqueness follows from the fact that any two solutions differ by a multiple of each $n_i$, hence by a multiple of $N$.
\end{proof}

\section{Primitive Roots}

\begin{definition}[Primitive Root]
Let $p$ be an odd prime. An integer $g$ is a \textit{primitive root modulo $p$} if the order of $g$ modulo $p$ is $p-1$, i.e., the smallest positive integer $k$ such that $g^k \equiv 1 \pmod{p}$ is $k = p-1$.
\end{definition}

Gauss proved that primitive roots exist for all prime moduli (Article 55). More generally, primitive roots exist modulo $n$ iff $n = 2, 4, p^k, 2p^k$ for an odd prime $p$.

\begin{algorithmthm}[Gauss's Algorithm for Discrete Logarithm]
Given a primitive root $g$ modulo $p$ and $a \not\equiv 0 \pmod{p}$, find $x$ such that $g^x \equiv a \pmod{p}$.

Gauss's method (Article 81) uses the fact that if $g$ is a primitive root, then every non-zero residue can be written uniquely as $g^x$. The algorithm reduces the problem by working with indices.
\end{algorithmthm}

\section{Python Implementation}
\label{sec:modular-python}

We implement modular arithmetic utilities and Gauss's algorithms in \texttt{python/modular.py}.

\begin{lstlisting}[style=python, caption={Modular arithmetic utilities (excerpt from python/modular.py)}]
# python/modular.py
"""Modular arithmetic utilities implementing Gauss's algorithms."""

from math import gcd
from typing import List, Tuple, Optional

def mod_inv(a: int, n: int) -> int:
    """Modular inverse using extended Euclidean algorithm.
    
    Returns x such that a*x == 1 (mod n).
    Raises ValueError if inverse doesn't exist.
    """
    def egcd(a: int, b: int) -> Tuple[int, int, int]:
        if b == 0:
            return (a, 1, 0)
        g, x1, y1 = egcd(b, a % b)
        return (g, y1, x1 - (a // b) * y1)
    
    g, x, _ = egcd(a, n)
    if g != 1:
        raise ValueError(f"No modular inverse for {a} mod {n}")
    return x % n

def chinese_remainder(residues: List[int], moduli: List[int]) -> int:
    """Solve system of congruences using Gauss's constructive proof.
    
    Args:
        residues: List of residues a_i
        moduli: List of pairwise coprime moduli n_i
        
    Returns:
        Solution x modulo product of moduli
    """
    if len(residues) != len(moduli):
        raise ValueError("Residues and moduli must have same length")
    
    # Verify pairwise coprime
    for i in range(len(moduli)):
        for j in range(i + 1, len(moduli)):
            if gcd(moduli[i], moduli[j]) != 1:
                raise ValueError(f"Moduli {moduli[i]} and {moduli[j]} are not coprime")
    
    N = 1
    for m in moduli:
        N *= m
    
    result = 0
    for a_i, n_i in zip(residues, moduli):
        N_i = N // n_i
        M_i = mod_inv(N_i, n_i)
        result += a_i * N_i * M_i
    
    return result % N

def primitive_root(p: int) -> int:
    """Find a primitive root modulo prime p.
    
    Uses Gauss's criterion: g is a primitive root iff
    g^((p-1)/q) != 1 (mod p) for all prime factors q of p-1.
    """
    if p == 2:
        return 1
    
    # Factor p-1
    def prime_factors(n: int) -> List[int]:
        factors = []
        d = 2
        while d * d <= n:
            while n % d == 0:
                factors.append(d)
                n //= d
            d += 1 if d == 2 else 2
        if n > 1:
            factors.append(n)
        return list(set(factors))
    
    factors = prime_factors(p - 1)
    
    for g in range(2, p):
        if all(pow(g, (p - 1) // q, p) != 1 for q in factors):
            return g
    
    raise ValueError(f"No primitive root found for {p}")

def discrete_log_gauss(g: int, a: int, p: int) -> Optional[int]:
    """Gauss's algorithm for discrete logarithm (baby-step giant-step).
    
    Finds x such that g^x == a (mod p), where g is a primitive root mod p.
    Returns None if no solution exists.
    """
    if a % p == 0:
        return None
    
    m = int((p - 1) ** 0.5) + 1
    
    # Baby steps: g^j for j = 0, 1, ..., m-1
    baby_steps = {}
    cur = 1
    for j in range(m):
        if cur not in baby_steps:
            baby_steps[cur] = j
        cur = (cur * g) % p
    
    # Giant steps: a * g^(-m*i)
    g_inv_m = pow(g, (p - 1) - m, p)  # g^(-m) mod p
    cur = a % p
    
    for i in range(m + 1):
        if cur in baby_steps:
            return i * m + baby_steps[cur]
        cur = (cur * g_inv_m) % p
    
    return None

def euler_totient(n: int) -> int:
    """Euler's totient function phi(n)."""
    result = n
    p = 2
    while p * p <= n:
        if n % p == 0:
            while n % p == 0:
                n //= p
            result -= result // p
        p += 1 if p == 2 else 2
    if n > 1:
        result -= result // n
    return result

def multiplicative_order(a: int, n: int) -> int:
    """Order of a modulo n (smallest k > 0 such that a^k == 1 mod n)."""
    if gcd(a, n) != 1:
        raise ValueError("a and n must be coprime")
    
    phi = euler_totient(n)
    # Check divisors of phi
    for d in range(1, phi + 1):
        if phi % d == 0 and pow(a, d, n) == 1:
            return d
    return phi
\end{lstlisting}

\section{Numerical Experiments}

\begin{example}[Chinese Remainder Theorem]
Solve:
\[
\begin{cases}
x \equiv 2 \pmod{3} \\
x \equiv 3 \pmod{5} \\
x \equiv 2 \pmod{7}
\end{cases}
\]

Using our implementation:
\begin{lstlisting}[style=python]
>>> from modular import chinese_remainder
>>> chinese_remainder([2, 3, 2], [3, 5, 7])
23
\end{lstlisting}

Verification: $23 \equiv 2 \pmod{3}$, $23 \equiv 3 \pmod{5}$, $23 \equiv 2 \pmod{7}$. The solution is unique modulo $105$.
\end{example}

\begin{example}[Primitive Roots]
Find primitive roots modulo $17$:
\begin{lstlisting}[style=python]
>>> from modular import primitive_root
>>> primitive_root(17)
3
\end{lstlisting}
The primitive roots modulo $17$ are $3, 5, 6, 7, 10, 11, 12, 14$ (all generators of $(\Z/17\Z)^\times$).
\end{example}

\begin{example}[Discrete Logarithm]
Find $x$ such that $3^x \equiv 13 \pmod{17}$:
\begin{lstlisting}[style=python]
>>> from modular import discrete_log_gauss
>>> discrete_log_gauss(3, 13, 17)
4
\end{lstlisting}
Indeed, $3^4 = 81 \equiv 13 \pmod{17}$.
\end{example}

\section{Exercises}

\begin{exercise}
Implement the Tonelli-Shanks algorithm for computing square roots modulo a prime. Compare with Gauss's method for $p \equiv 3 \pmod{4}$.
\end{exercise}

\begin{exercise}
Write a function to compute the index (discrete logarithm) table for a given primitive root modulo $p$. Use it to solve $g^x \equiv a \pmod{p}$ by table lookup.
\end{exercise}

\begin{exercise}
Prove that if $g$ is a primitive root modulo $p$, then $g^k$ is also a primitive root iff $\gcd(k, p-1) = 1$. How many primitive roots exist modulo $p$?
\end{exercise}

\begin{exercise}
Implement the Pohlig-Hellman algorithm for discrete logarithms when $p-1$ is smooth. Compare performance with baby-step giant-step.
\end{exercise}

\section{Historical Notes}

Gauss's \textit{Disquisitiones Arithmeticae} was written when he was just 21 years old (published at 24). The work introduced the notation $a \equiv b \pmod{n}$ and systematically developed the theory of congruences. Article 36--44 contains the Chinese Remainder Theorem with Gauss's own proof. Articles 52--60 establish the existence of primitive roots for prime moduli. The discrete logarithm (index) calculus in Articles 81--85 was used by Gauss for extensive numerical computations.

\section{References}

\begin{itemize}
    \item \cite{gauss-disq} -- \textit{Disquisitiones Arithmeticae}, Articles 1--85
    \item \cite{edwards-fermat} -- Edwards, H.M., \textit{Fermat's Last Theorem: A Genetic Introduction to Algebraic Number Theory}
    \item \cite{stillwell-number-theory} -- Stillwell, J., \textit{Elements of Number Theory}
    \item \cite{cohen-number-theory} -- Cohen, H., \textit{Number Theory: Volume I: Tools and Diophantine Equations}
\end{itemize}